\section{Conclusions}
Intelligent dynamic pricing strategy under the background of big data analysis is indispensable in shipping container industry. Our team has drawn the following conclusions when designing pricing strategy based on existing data.
\begin{itemize}
    \item We select ten flow directions based on historical revenue, order quantity, freight rate, and throughput. The results are presented in Table \ref{tbl:Mock-TOPSIS Evaluation Results.}.These ten samples are proved to be meaningful and helpful in the research process.
    \item The time distribution of customers in different flow directions for ordering containers is diverse. The reason behind this is that customers will consider buying according to the rise and fall of the freight rate, thus reflecting the negative correlation between sales volume and shipping rate.
    \item Fluctuations in market demand will affect sales volume and further influence changes in freight rates.
    \item The current pricing strategy is beneficial to increase revenue, but it has limitations due to insufficient considerations.
    \item We design a pricing system to determine whether to adjust the price and the extent of the price adjustment based on the threshold, and test its rationality, and achieved very satisfactory results.
\end{itemize}

However, due to the lack of data and the pressing time, we made many bold assumptions before conducting this research. Furthermore, our future work can be realized when conditions permit.
\begin{itemize}
    \item If we can get the order statistics of the contract market for a long time, we can know the container sales volume generated by the contract market in the short term, and predict the sales volume in the spot market more accurately when the inventory constraints are known\parencite{dynamic-pricing}.
    \item It is more realistic and valuable to maximize the benefits under the consideration of costs. If we can know the composition of container transportation costs, this step can be achieved by us.
\end{itemize}